\begin{itemize}
    \item We conduct additional experiments using extrinsic local regression
    for manifold-valued responses and Euclidean covariates
    \citep{linExtrinsicLocalRegression2017}. The observed covariate should not
    be the latent projection index used by PME: if that smooth ordering were
    observed, one could directly regress the two ambient coordinates on the
    ordering. We therefore generate a Euclidean covariate once, condition on
    the resulting design, and relate it to the flower phase through a
    deterministic, nonmonotone map. All stochastic variation after fixing the
    design is introduced at the response level.

    \item For each $i=1,\ldots,N$, generate the fixed design once as
    \begin{align}
        t_i &\overset{\mathrm{iid}}{\sim} \operatorname{Uniform}(-1,1),\\
        u_i &= \left\{\frac{1+\sin(\pi t_i)}{2}\right\}\bmod 1.
    \end{align}
    After this initial generation, inference is conditional on
    $\{t_i\}_{i=1}^N$. In particular, there is no stochastic error in either
    $t_i$ or $u_i$. Because $t\mapsto u(t)$ is nonmonotone, $t_i$ is not the
    latent ordering or projection index of the flower.

    Define the flower embedding $F:S^1\rightarrow\mathbb{R}^2$ by
    \begin{align}
        \theta(u) &= 2\pi u,\\
        r(u) &= 1+a\cos\{5\theta(u)\},\\
        F(u) &= r(u)
        \begin{pmatrix}
            \cos\{\theta(u)\}\\
            \sin\{\theta(u)\}
        \end{pmatrix},
    \end{align}
    with flower amplitude $a=0.30$. The conditional mean response is
    \[
        \mu_i=\mu(t_i)=F(u_i),
    \]
    and the observed ambient response is
    \begin{align}
        X_i &= \mu_i+\varepsilon_i,\\
        \varepsilon_i &\overset{\mathrm{iid}}{\sim}
        \mathcal{N}_2\!\left(
        \begin{pmatrix}0\\0\end{pmatrix},
        \sigma_X^2 I_2\right),
    \end{align}
    where $\sigma_X=0.04$. Thus, unlike the previous phase-variability
    experiment, the present experiment has no latent phase error; the only
    noise is additive response noise in the ambient space.

    \item For a candidate PME smoothing parameter $\lambda_j$ and a given
    training fold, we proceed as follows:
    \begin{enumerate}
        \item Fit PME to the training responses only, obtaining the manifold
        estimate $\widehat S_{j}$.
        \item Project each training response onto $\widehat S_j$ to obtain
        \[
            \widetilde Y_{i,j}
            =P_{\widehat S_j}(X_i).
        \]
        \item Form the training pairs $(t_i,\widetilde Y_{i,j})$ and fit both
        extrinsic regression estimators proposed by
        \citet{linExtrinsicLocalRegression2017}.

        For extrinsic kernel regression, the ambient estimate is
        \[
            \widehat m_{K,j}(t)
            =\frac{\sum_iK_h(t_i-t)\widetilde Y_{i,j}}
            {\sum_iK_h(t_i-t)},
        \]
        followed by
        \[
            \widehat\mu_{E,K,j}(t)
            =P_{\widehat S_j}\{\widehat m_{K,j}(t)\}.
        \]
        For extrinsic local-linear regression, we calculate
        \begin{align}
            (\widehat\beta_0,\widehat\beta_1)
            =\arg\min_{\beta_0,\beta_1}
            \sum_i K_h(t_i-t)
            \left\|\widetilde Y_{i,j}
            -\beta_0-\beta_1(t_i-t)\right\|^2,
        \end{align}
        and then set
        \[
            \widehat\mu_{E,LL,j}(t)
            =P_{\widehat S_j}\{\widehat\beta_0(t)\}.
        \]
    \end{enumerate}

    \item We use nested five-fold cross-validation. For each candidate
    $\lambda$ and outer fold, PME is fitted to the training responses and the
    regression bandwidth is selected by leave-one-out prediction of the
    projected training responses. The resulting pipeline is evaluated against
    the raw, unprojected outer-fold responses, which are used only to select
    $\lambda$. RMSE against the known conditional mean $\mu(t_i)$ is reported
    only as a simulation-oracle diagnostic and is not used for tuning.

    The candidate grid was
    \[
        \lambda\in\{10^{-14},10^{-13},\ldots,10^{-4}\},
    \]
    with regression bandwidths ranging from $0.004$ to $0.38$. Both procedures
    selected $\lambda=10^{-8}$. For extrinsic kernel regression, the mean
    fold-selected bandwidth was $0.0064$, the full-data bandwidth was $0.006$,
    and the raw and simulation-oracle RMSEs were $0.06626$ and $0.03483$.
    For extrinsic local-linear regression, the corresponding values were
    $0.0092$, $0.008$, $0.06044$, and $0.02215$. All 55 fold-specific PME fits
    converged.
\end{itemize}

\begin{figure}[H]
    \centering
    \includegraphics[width=1\linewidth]{08_2026_fig/kfold_performance.png}
    \caption{Five-fold out-of-fold performance for extrinsic kernel regression
    (left) and extrinsic local-linear regression (right). At each value of
    $\lambda$ and outer fold, the regression bandwidth is chosen by
    leave-one-out prediction of the projected training responses. The orange
    curve is mean RMSE against the untouched noisy outer-fold responses and is
    used to select $\lambda$. The blue dashed curve is RMSE against the known
    conditional mean, evaluated at the same fold-selected bandwidths, and is
    shown only as an oracle diagnostic. Both methods select
    $\lambda=10^{-8}$.}
    \label{fig:kfold_performance}
\end{figure}

\begin{figure}[H]
    \centering
    \includegraphics[width=1\linewidth]{08_2026_fig/kfold_fits.png}
    \caption{Fixed-design responses with response-noise standard deviation
    $\sigma_X=0.04$ (left) and the full-data PME fit at the cross-validation-
    selected value $\lambda=10^{-8}$ (right). The black dashed curve is the
    true flower and the blue curve is the selected PME manifold. The kernel
    and local-linear downstream criteria selected the same PME smoothing
    parameter and therefore yield the same manifold fit.}
    \label{fig:kfold_fits}
\end{figure}
